\documentclass[12pt,letterpaper]{article}
\usepackage[T1]{fontenc}
\usepackage{microtype,amsmath,amssymb,amsthm,mathtools}
\usepackage{newtxtext,newtxmath}
\usepackage[letterpaper,margin=1in]{geometry}
\usepackage[hidelinks]{hyperref}
\usepackage{enumitem}
\usepackage{booktabs}
\usepackage{tikz}
\newtheorem{theorem}{Theorem}[section]
\newtheorem{lemma}[theorem]{Lemma}
\newtheorem{proposition}[theorem]{Proposition}
\newtheorem{corollary}[theorem]{Corollary}
\theoremstyle{remark}
\newtheorem{remark}[theorem]{Remark}
\newtheorem{example}[theorem]{Example}
\numberwithin{equation}{section}
\newcommand{\C}{\mathbb C}
\newcommand{\Hh}{\mathcal H}
\newcommand{\D}{\mathcal D}
\newcommand{\Ker}{\operatorname{Ker}}
\newcommand{\ran}{\operatorname{Ran}}
\newcommand{\Dom}{\operatorname{Dom}}
\newcommand{\spn}{\operatorname{span}}
\newcommand{\Hp}{H^2(\C_+)}
\newcommand{\Nplus}{N^+(\C_+)}
\newcommand{\dd}{\,\mathrm d}
\title{Spectral Invariance under Grosswald--Schnitzer Prime Perturbations\\[0.4em]
\large Draft 2\\[0.25em]
\normalsize Bounded comparison operators, unbounded counterexamples, and a common Hardy-space model}
\author{David Joyner\\\texttt{wdjoyner@gmail.com}}
\date{November 28, 2026}
\begin{document}
\maketitle
\begin{abstract}
This paper was prepared with assistance from ChatGPT (OpenAI) and Claude (Anthropic); the specific contributions are described in the Acknowledgments and the revision log.
We investigate whether the zero-preserving perturbations of primes introduced by Grosswald and Schnitzer in 1978 induce boundedly invertible comparisons between the zeta-associated dilation operators of D. Joyner. We prove an elementary sublinear growth estimate for the logarithm of the Euler-product quotient $Q_A=\zeta_A/\zeta$ in every right half-plane $\Re s\ge\delta>0$, and deduce that $Q_A$ and its reciprocal belong to the Smirnov class and that $Q_A$ is outer. Applied to the Hardy-space representation of Joyner's weighted Sobolev space, this proves equality of the \emph{closed joint kernels}, not just their eigenvalue multisets, for all Grosswald--Schnitzer perturbations. We distinguish this equality from similarity of the full operators. We give sufficient conditions for genuinely boundedly invertible comparison operators, as well as arbitrarily small perturbations for which $Q_A$ and $1/Q_A$ are both unbounded on the relevant half-plane. Unconditionally, the closed joint kernel on the space $S_1$ is $\{0\}$ if and only if the Riemann hypothesis holds. Taking the November 2026 certified formal proof of RH as an external input, the closed joint kernel is therefore $\{0\}$ on $S_1$, and on $S_k$ with $k\ge2$ it is the closed span of the critical-line zero powers and their logarithmic companions. The results use standard Hardy-space facts and require independent novelty assessment before publication; this paper does not itself prove or verify RH.
\end{abstract}

\clearpage
\setcounter{tocdepth}{2}
\tableofcontents
\clearpage

\section{Introduction}
\subsection{Motivation and principal results}
The Grosswald--Schnitzer theorem says that moving each prime within
its prime gap does not move the zeros of the associated zeta function
in $\Re s>0$.  This paper asks how much of that invariance survives
at the level of the zeta-associated dilation operators of the
author's 1981 manuscript (see the corrected edition in the
\texttt{zeta-operators-graph} folder of this repository).

We distinguish invariance of the zero-associated common kernel
from bounded left equivalence or similarity of full summation operators.
The first is complete: Theorem~\ref{thm:main} shows that the
\emph{closed} joint kernel does not depend on the perturbation.
The second depends on the perturbation in a delicate way
(Theorem~\ref{thm:compare}); its exact classification for
endpoint shifts is given in the companion paper \cite{PP26}.
The proof uses Hardy-space and Smirnov-class factorization: the
Euler quotient $Q_A$ may be unbounded, but it is always an
\emph{outer} function, so it introduces no new inner factor.
Combined with RH, Theorem~\ref{thm:main} determines the closed joint
kernel completely (Corollary~\ref{cor:rh-kernel}).

The simplest example to keep in mind is the endpoint perturbation
$a_n=p_{n+1}$, which moves every prime to its successor.  Its Euler
quotient telescopes to $Q_A(s)=1-2^{-s}$, which is bounded and
bounded away from zero on every half-plane $\Re s\ge\delta>0$.
\subsection{Literature review}
Grosswald and Schnitzer introduced the zero-preserving prime-gap
perturbations \cite{GS}.  Beurling's invariant-subspace
theory is a precursor of the Hardy-space model \cite{Beurling}.
The Smirnov-class material relates to work of Choa--Kim--Shapiro
\cite{SmirnovReference} and Fricain--Hartmann--Ross--Timotin
\cite{Fricain}.  The prior spectral questions were formulated
in the accompanying working note \cite{JoynerNote}.
The companion paper \cite{PP26} classifies bounded comparison
for endpoint perturbations and for perturbations with
rationally independent frequencies.

\paragraph{The November 2026 RH formalization.}
On November 20, 2026, Levent Alp\"oge released a Lean formalization
of a proof of RH, found with Anthropic's Claude \cite{Ru26}.
The Lean FRO and Mathlib Initiative certified
the statement and formal proof on November 21 \cite{LF26}.
Rubinstein records independent checks of the original formalization
and its shorter successor \cite{Ru26}. We use this development as an
external input in Section~\ref{sec:rh-2026}, not as a
replacement for the closure arguments of the Hardy-space model.
For $k\ge2$, RH locates the
nontrivial zero powers strictly within the admissible exponent
half-plane of $S_k$; for $k=1$, they lie on its boundary.
Together with Theorem~\ref{thm:main}, this determines the closed
joint kernel for every $k$ (Corollary~\ref{cor:rh-kernel}).
RH proves neither zero simplicity nor similarity of the full
operator pairs. Human-readable expositions remain under development,
and the present paper has not independently audited the formal proof.
\subsection{Notation and conventions}
Write $A=(a_n)$ for a prime-gap perturbation and
$Q_A=\zeta_A/\zeta$ for its Euler quotient.
Fix $k\ge1$ and $\delta=1/(2k)$.
The space $S_k$ is weighted Sobolev,
$H^2(\mathbb C_+)$ is the Hardy model,
and $N^+(\mathbb C_+)$ is the Smirnov class.
The dilation generator is $\Delta=y\,d/dy$, so that
$\Delta y^s=sy^s$, and $D_cf(y)=f(y/c)$ denotes dilation by $c>1$.
The operators $T_A^{(r)}$ are initially analytic-symbol
operators; literal convergence in the critical strip is not assumed.

\section{Setting and main conclusions}
Let $p_1=2<p_2=3<\cdots$ be the primes. Choose real numbers
\begin{equation}\label{eq:GS}
 p_n\le a_n\le p_{n+1}\quad(n\ge1),
\end{equation}
and define for $\Re s>1$
\begin{equation}\label{eq:Q}
 \zeta_A(s)=\prod_{n\ge1}(1-a_n^{-s})^{-1},\qquad
 Q_A(s)=\frac{\zeta_A(s)}{\zeta(s)}
 =\prod_{n\ge1}\frac{1-p_n^{-s}}{1-a_n^{-s}}.
\end{equation}
Grosswald--Schnitzer \cite{GS} proves that $Q_A$ extends holomorphically without zeros to $\Re s>0$, hence $\zeta_A=Q_A\zeta$ there. We will also derive the required convergence directly below.

Fix an integer $k\ge1$, put $\delta=1/(2k)$, and work in the \emph{full} weighted Sobolev space
\begin{equation}\label{eq:Sk}
 \|f\|_{S_k}^2=\sum_{j=0}^k\int_0^1|f^{(j)}(y)|^2y^{2j-1-1/k}\dd y,
\end{equation}
with no added boundary condition at $y=1$. Its standard algebraic domain is
\[
 \mathcal D_k=\spn\{y^s(\log y)^j:\Re s>\delta,\ j\ge0\}.
\]
A power $y^s$ lies in $S_k$ exactly when $\Re s>\delta$: each term
of \eqref{eq:Sk} has integrand $|s(s-1)\cdots(s-j+1)|^2y^{2\Re s-1-1/k}$.

For $r=2,3$, define the holomorphic symbols on $\Re s>\delta$
\begin{equation}\label{eq:symbols}
 m_r(s)=(1-r^{1-s})\zeta(s),\qquad m_{A,r}(s)=Q_A(s)m_r(s).
\end{equation}
Here $m_r$ is entire: the zero of $1-r^{1-s}$ at $s=1$ cancels the pole of $\zeta$. Define operators on $\mathcal D_k$ by
\begin{equation}\label{eq:operators}
 T^{(r)}[y^s(\log y)^j]=\partial_s^j(m_r(s)y^s),\quad
 T_A^{(r)}[y^s(\log y)^j]=\partial_s^j(m_{A,r}(s)y^s).
\end{equation}
These definitions agree with the original alternating dilation sums for $T^{(r)}$. For $T_A^{(r)}$, they are \emph{analytic-symbol definitions}; no assertion of convergence of a literal generalized-prime dilation series for $\Re s\le1$ is implicit.

\begin{theorem}[Spectral invariance]\label{thm:main}
For every choice \eqref{eq:GS} and every $k\ge1$, the four operators in \eqref{eq:operators} are densely defined and closable on $S_k$, and
\begin{equation}\label{eq:main}
 \boxed{\Ker\overline{T_A^{(2)}}\cap\Ker\overline{T_A^{(3)}}
 =\Ker\overline{T^{(2)}}\cap\Ker\overline{T^{(3)}}.}
\end{equation}
Both sides equal the closed $S_k$-span of
\[
 y^\rho(\log y)^j,\qquad
 \zeta(\rho)=0,\quad \Re\rho>\delta,\quad 0\le j<\operatorname{ord}_\rho\zeta.
\]
Thus the two zero-associated spectral models are literally the same closed Hilbert subspace, with the same induced norm, and the identity is a unitary identification intertwining their restrictions of the dilation generator with the same chosen domain.
\end{theorem}

\begin{theorem}[Bounded comparison versus arbitrary perturbations]\label{thm:compare}
The diagonal map $C_{Q_A}y^s=Q_A(s)y^s$ extends to a boundedly invertible operator on $S_k$ if and only if
\[
 Q_A,\,1/Q_A\in H^\infty(\Re s>\delta).
\]
There exist nontrivial \emph{infinite} perturbations satisfying this condition. Conversely, there exist perturbations with $0<a_n-p_n<2^{-n}$ for all $n$ such that \emph{both} $Q_A$ and $1/Q_A$ are unbounded even on $\Re s>1$. These unbounded examples still satisfy Theorem~\ref{thm:main}.
\end{theorem}

\begin{example}[Moving only the first prime]\label{ex:first-prime}
Take $a_1=5/2$ and $a_n=p_n$ for $n\ge2$.  Then
\[
 Q_A(s)=\frac{1-2^{-s}}{1-(5/2)^{-s}},\qquad
 m_{A,r}(s)=\frac{1-2^{-s}}{1-(5/2)^{-s}}\,(1-r^{1-s})\zeta(s).
\]
For example $Q_A(2)=\tfrac{3/4}{21/25}=\tfrac{25}{28}$, so the
eigenvalue of $T_A^{(r)}$ on $y^2$ is $25/28$ times that of $T^{(r)}$.
On $\Re s>\delta$ both $|1-2^{-s}|$ and $|1-(5/2)^{-s}|$ lie between
$1-2^{-\delta}$ and $1+2^{-\delta}$, so $Q_A$ and $1/Q_A$ are bounded
there.  The full operators therefore differ, but by
Theorem~\ref{thm:compare} only by a boundedly invertible diagonal
factor, and by Theorem~\ref{thm:main} their joint kernels coincide.
\end{example}

\begin{remark}[The order of the Sobolev norm]\label{rem:order}
The corrected edition of the author's 1981 manuscript, in the
companion folder of this repository, uses the same weights but sums the derivatives up to a fixed order $N\ge2$,
independent of $k$.  The admissible half-plane $\Re s>\delta$ is the
same.  Every argument below goes through for that space with $k$
replaced by $N$ in the exponent of $(\Delta+c)^k$ in \eqref{eq:Uk};
then $U(y^s)(z)=(s+c)^N/(z+s-\delta)$.  In particular,
Theorem~\ref{thm:main} and Corollary~\ref{cor:rh-kernel} hold for
the spaces of the 1981 edition as well.
\end{remark}

\section{A uniform sublinear growth estimate}
We first prove the analytic estimate which makes the invariant-space argument work.

\begin{lemma}[Growth of the comparison quotient]\label{lem:growth}
For each $0<\delta<1$, one may choose an analytic logarithm $L_A=\log Q_A$ on $\Re s>0$ such that
\begin{equation}\label{eq:sublinear}
 \left|L_A(s)\right|\le C_\delta(1+|s|)^{1-\delta},\qquad \Re s\ge\delta,
\end{equation}
where the constant can be chosen \emph{uniformly in all sequences} satisfying \eqref{eq:GS}.
\end{lemma}
\begin{proof}
For $x\ge2$ and $\Re s>0$ use the branch
$h_s(x)=\log(1-x^{-s})=-\sum_{m\ge1}x^{-ms}/m$.
For $\Re s\ge\delta$ one has
\[
 |h_s(x)|\le C_\delta x^{-\delta},\qquad
 |h_s'(x)|=\left|\frac{s x^{-s-1}}{1-x^{-s}}\right|
 \le C_\delta |s|x^{-1-\delta}.
\]
Writing $g_n=p_{n+1}-p_n$, it follows that
\begin{equation}\label{eq:each}
 |h_s(p_n)-h_s(a_n)|
 \le C_\delta\min\big\{p_n^{-\delta},\ |s|g_n p_n^{-1-\delta}\big\}.
\end{equation}
The series of logarithmic differences is absolutely and locally uniformly convergent in $\Re s>0$: for $s$ in a compact set, use the second bound and the convergence of $\sum_n g_np_n^{-1-\epsilon}$ for each $\epsilon>0$. The latter follows by comparison with the integral of $x^{-1-\epsilon}$ over the disjoint prime intervals $[p_n,p_{n+1}]$, since Bertrand's postulate gives $p_{n+1}<2p_n$.

Take $R=2+|s|$. Splitting \eqref{eq:each} at $p_n=R$ gives
\[
 |L_A(s)|\le C_\delta\left(\sum_{p_n\le R}p_n^{-\delta}
 +|s|\sum_{p_n>R}g_np_n^{-1-\delta}\right).
\]
The first sum is $O_\delta(R^{1-\delta})$ by comparison with all integers. The second sum is $O_\delta(|s|R^{-\delta})$, again by the disjoint intervals and Bertrand's postulate. This yields \eqref{eq:sublinear}.
\end{proof}

\begin{remark}[Why primewise closeness does not give an $H^\infty$ bound]
An estimate based on the mean-value theorem has a factor $|s|$, whereas the trivial estimate does not measure $|a_n-p_n|$. Splitting the sum at $p_n\asymp|s|$ produces the sublinear growth exponent $1-\delta$, but it \emph{does not} produce a uniform bound as $|\Im s|\to\infty$. In particular, control on every compact vertical strip is weaker than an $H^\infty$ estimate on an entire half-plane.
\end{remark}

\section{The Smirnov-class mechanism}
Let $\C_+=\{z:\Re z>0\}$ and $\Nplus$ denote the Smirnov class: functions expressible as $b/a$ with $a,b\in H^\infty(\C_+)$ and $a$ outer. We use the standard facts that products of outer functions are outer, the outer functions in $H^\infty$ have dense multiples in $H^2$, and a member of $N^+$ is invertible in $N^+$ precisely when it is outer \cite{SmirnovReference}.

\begin{lemma}[Invertibility in $N^+$]\label{lem:outer}
The functions $z\mapsto Q_A(z+\delta)$ and $z\mapsto Q_A(z+\delta)^{-1}$ belong to $\Nplus$, and $Q_A(z+\delta)$ is outer there.
\end{lemma}
\begin{proof}
Choose $\beta$ with $1-\delta<\beta<1$ and let
\[
 E_\beta(z)=\exp\big[-(1+z)^\beta\big],\qquad\Re z>0,
\]
using the principal branch. This is a bounded, nonvanishing \emph{outer} function: its logarithmic boundary modulus is integrable with weight $(1+t^2)^{-1}$, it extends analytically across each finite boundary point, and $-\log|E_\beta(x)|=o(x)$ as $x\to+\infty$, excluding a singular inner factor at infinity. Since
$\Re(1+z)^\beta\ge \cos(\pi\beta/2)|1+z|^\beta$ for $\Re z\ge0$, Lemma~\ref{lem:growth} gives
\[
 E_\beta(z)Q_A(z+\delta),\qquad
 E_\beta(z)/Q_A(z+\delta)\in H^\infty(\C_+).
\]
Thus both $Q_A$ and its reciprocal are quotients of bounded analytic functions by the outer denominator $E_\beta$. Since a product of two Smirnov functions is again Smirnov and their product is $1$, their canonical inner factors must be constant. Hence $Q_A$ is outer.
\end{proof}

\begin{lemma}[Ranges of unbounded Hardy multipliers]\label{lem:range}
Suppose $0\ne f\in\Nplus$ and write its canonical factorization $f=I_fO_f$ (inner times outer). The maximal multiplication operator
\[
 M_f:F\mapsto fF,\quad\Dom M_f=\{F\in H^2(\C_+):fF\in H^2(\C_+)\}
\]
is closed and densely defined, and
\begin{equation}\label{eq:range}
 \overline{\ran M_f}=I_fH^2(\C_+).
\end{equation}
Moreover the span of all Hardy reproducing kernels $K_w(z)=(z+\overline w)^{-1}$ is a graph core for $M_f^*$.
\end{lemma}
\begin{proof}
Write $f=b/a$ with $a,b\in H^\infty$ and $a$ outer. Since $aH^2\subset\Dom M_f$ and $aH^2$ is dense, $M_f$ is densely defined. It is closed because convergence in $H^2$ implies locally uniform convergence on $\C_+$, and multiplication by a holomorphic function respects such convergence. Also $M_f(aF)=bF$, so $\overline{\ran M_f}\supset\overline{bH^2}=I_bH^2=I_fH^2$.
For the reverse inclusion, if $fF=G\in H^2$ then $bF=aG$. The inner factor $I_b$ divides the inner factor of $aG$, and $a$ is outer, so $I_b$ divides $G$. This proves \eqref{eq:range}.

For $w\in\C_+$, bounded point evaluation gives $M_f^*K_w=\overline{f(w)}K_w$ (with inner product linear in the first variable). If a pair $(F,G)\in H^2\oplus H^2$ is orthogonal to every graph vector $(K_w,M_f^*K_w)$, evaluating at all $w$ shows $F=-fG$. Thus $G\in\Dom M_f$ and $F=-M_fG$: the orthogonal complement of these kernel graph vectors coincides with that of the full graph of $M_f^*$. Their span is therefore a graph core.
\end{proof}

\begin{example}[Outer, bounded and unbounded]\label{ex:outer}
For the endpoint shift $a_n=p_{n+1}$ one has
$Q_A(z+\delta)=1-2^{-z-\delta}$.  On $\C_+$ its modulus lies between
$1-2^{-\delta}$ and $1+2^{-\delta}$.  A function that is bounded and
bounded away from zero is invertible in $H^\infty(\C_+)$, and an
invertible element of $H^\infty$ is outer, so here
Lemma~\ref{lem:outer} is immediate.  The lemma has real content for
the perturbations of Proposition~\ref{prop:counter} below.  There
$Q_A(z+\delta)$ is outer but \emph{not} bounded, so it is not a bounded
multiplier, yet it still contributes no inner factor.
\end{example}

\section{Proof of spectral invariance}\label{sec:proof}
Let $U_k:S_k\to H^2(\C_+)$ denote the bounded isomorphism of the earlier spectral-synthesis note \cite{JoynerNote}. With $\Delta=y\,d/dy$ and $c>0$ it is the Laplace transform of $(\Delta+c)^kf$ after the logarithmic change $t=-\log y$ and the weight $e^{\delta t}$:
\[
 (U_kf)(z)=\int_0^\infty e^{-zt}\,e^{\delta t}\,\big((\Delta+c)^kf\big)(e^{-t})\dd t .
\]
For the reader's convenience, here is why $U_k$ is an isomorphism.
Since $dy/y=dt$, the norm \eqref{eq:Sk} equals
$\sum_{j\le k}\|e^{\delta t}(y^jD^jf)(e^{-t})\|_{L^2(0,\infty)}^2$.
Each $y^jD^j$ is a polynomial in $\Delta$ of degree $j$, and conversely,
so $\|f\|_{S_k}$ is equivalent to
$\sum_{i\le k}\|e^{\delta t}(\Delta^if)(e^{-t})\|_{L^2}$.
In the variable $t$, $\Delta=-d/dt$, and conjugation by $e^{\delta t}$
turns $\Delta+c$ into $-d/dt+(c+\delta)$.  The latter is a bijection from
the Sobolev space $W^{1,2}(0,\infty)$ onto $L^2(0,\infty)$.  Its inverse is
convolution with the integrable kernel $e^{-(c+\delta)u}\mathbf 1_{u\ge0}$,
taken from the right, so no boundary condition at $t=0$ ($y=1$) is
imposed.  Hence $f\mapsto e^{\delta t}((\Delta+c)^kf)(e^{-t})$ is an
isomorphism of $S_k$ onto $L^2(0,\infty)$, and the Paley--Wiener theorem
identifies the Laplace image of $L^2(0,\infty)$ with $H^2(\C_+)$.
The essential property of $U_k$ is
\begin{equation}\label{eq:Uk}
 U_k(y^s)(z)=\frac{(s+c)^k}{z+s-\delta},\qquad \Re s>\delta.
\end{equation}
Thus $\mathcal D_k$ maps to the span of the Hardy evaluation kernels and their parameter derivatives. Any real-symmetric holomorphic symbol $f(s)$ in the shifted Smirnov class determines, by \eqref{eq:operators}, a closable algebraic operator whose closure is
\begin{equation}\label{eq:adjointmodel}
 U_k^{-1}M_{f(z+\delta)}^*U_k,
\end{equation}
by Lemma~\ref{lem:range} and the reproducing-kernel graph core. This includes our $m_r$ and $m_{A,r}$.

\begin{proof}[Proof of Theorem~\ref{thm:main}]
The functions $m_r(z+\delta)$ are in $\Nplus$: they have polynomial growth in the closed half-plane, so division by a sufficiently high power of $1+z$ gives a bounded numerator and a bounded outer denominator. Lemma~\ref{lem:outer} implies that $m_{A,r}(z+\delta)=Q_A(z+\delta)m_r(z+\delta)$ has \emph{exactly the same inner factor} as $m_r(z+\delta)$.

The only zeros in $\Re s>\delta$ of $m_r$ are the zeta zeros and the additional points
\[
 s=1+\frac{2\pi i n}{\log r},\qquad n\in\mathbb Z\setminus\{0\}.
\]
There is no additional common zero for $r=2,3$, because $\log2/\log3$ is irrational. Also the $m_r$ are entire of polynomial half-plane growth, so there are no finite-boundary singular inner factors; their limit $m_r(\sigma)\to1$ as $\sigma\to+\infty$ excludes a singular inner factor at infinity. Consequently the greatest common inner divisor $B_\delta$ of the two symbols is just the Blaschke product of the zeta zeros $\rho$ with $\Re\rho>\delta$, repeated according to multiplicity.

For any $f_1,f_2\in\Nplus$, Lemma~\ref{lem:range} and Beurling's theorem give
\[
 \overline{\ran M_{f_1}+\ran M_{f_2}}
 =\gcd(I_{f_1},I_{f_2})H^2(\C_+).
\]
Taking orthogonal complements and using \eqref{eq:adjointmodel} proves
\[
 U_k\left(\Ker\overline{T_A^{(2)}}\cap\Ker\overline{T_A^{(3)}}\right)
 =H^2\ominus B_\delta H^2
 =U_k\left(\Ker\overline{T^{(2)}}\cap\Ker\overline{T^{(3)}}\right).
\]
Completeness of the reproducing kernels and derivative kernels at the zeros of $B_\delta$ yields the asserted spectral synthesis, exactly as in the earlier note \cite{JoynerNote}. Since the two closed subspaces are equal, their inherited Hilbert norms are equal; if the dilation generator is equipped with the same restriction/domain in both models, the identity intertwines it.
\end{proof}

\begin{example}[The first zeta zero]\label{ex:first-zero}
Let $\rho_1=\tfrac12+i\gamma_1$ with $\gamma_1\approx14.1347$ be the
first nontrivial zero of $\zeta$.  For $k=2$ we have $\delta=\tfrac14<\Re\rho_1$,
so $y^{\rho_1}\in S_2$.  Under \eqref{eq:Uk} it maps to a multiple of the
Hardy kernel $K_w$ with $w=\overline{\rho_1}-\tfrac14=\tfrac14-i\gamma_1$.
That kernel lies in $H^2\ominus B_{1/4}H^2$, since $\langle Bh,K_w\rangle$ is a
constant multiple of $B(w)h(w)$, which vanishes when $B(w)=0$.  Here $B(w)=0$ because $\overline{\rho_1}=\tfrac12-i\gamma_1$ is
also a zero of $\zeta$.  Thus $y^{\rho_1}$ lies in the joint kernel of
$\overline{T_A^{(2)}}$ and $\overline{T_A^{(3)}}$ for \emph{every}
Grosswald--Schnitzer perturbation $A$.  For $k=1$, $\delta=\tfrac12$ and
$y^{\rho_1}\notin S_1$.
\end{example}

\section{When comparisons really are bounded}\label{sec:bounded}
\begin{proposition}[An $H^\infty$ criterion]\label{prop:Hinfty}
For any holomorphic, real-symmetric $Q$ on $\Re s>\delta$, the diagonal map $C_Qy^s=Q(s)y^s$ extends boundedly to $S_k$ if and only if $Q\in H^\infty(\Re s>\delta)$. It is boundedly invertible if and only if both $Q$ and $1/Q$ belong to $H^\infty$ there.
\end{proposition}
\begin{proof}
Necessity follows from evaluating the norm on each nonzero vector $y^s$. Sufficiency follows from the bounded Hardy multiplication adjoint conjugated by $U_k$. The inverse, if bounded, must act as multiplication by $1/Q$ on the total family of power functions.
\end{proof}

\begin{proposition}[Sparse infinite perturbations]\label{prop:sparse}
If $I=\{n:a_n\ne p_n\}$ and
\begin{equation}\label{eq:sparse}
 \sum_{n\in I}p_n^{-\delta}<\infty,
\end{equation}
then $Q_A,1/Q_A\in H^\infty(\Re s>\delta)$.
\end{proposition}
\begin{proof}
For $\Re s>\delta$ and any changed index,
\[
 \left|\log\frac{1-p_n^{-s}}{1-a_n^{-s}}\right|
 \le C_\delta(p_n^{-\delta}+a_n^{-\delta})
 \le2C_\delta p_n^{-\delta}.
\]
The uniformly convergent logarithm therefore has bounded real part, above and below. Exponentiating yields the two bounds.
\end{proof}

\begin{remark}[Concrete examples]
Replacing only $2$ by $5/2$ produces
\[
 Q_A(s)=\frac{1-2^{-s}}{1-(5/2)^{-s}},
\]
which is bounded and bounded away from zero for $\Re s>\delta$. For a genuinely infinite example, select indices $n_j$ so that $p_{n_j}>2^{j^2}$, let $a_{n_j}=(p_{n_j}+p_{n_j+1})/2$, and leave all other primes unchanged. Condition \eqref{eq:sparse} holds for every fixed $\delta>0$.
\end{remark}

\section{Counterexample: arbitrarily small shifts can destroy boundedness}\label{sec:counter}
\begin{proposition}[Generic non-$H^\infty$ comparison]\label{prop:counter}
There are sequences satisfying \eqref{eq:GS}, with $0<a_n-p_n<2^{-n}$ for every $n$, for which
\[
 Q_A\notin H^\infty(\Re s>1),\qquad
 1/Q_A\notin H^\infty(\Re s>1).
\]
Consequently neither comparison is bounded on $S_k$ for any $k\ge1$.
\end{proposition}
\begin{proof}
Choose $a_n$ successively in
$(p_n,\min\{p_n+2^{-n},p_{n+1}\})$, avoiding the countable set for which $\log a_n$ lies in the rational linear span of all $\log p_j$ and the earlier $\log a_j$. Then the family
\[\{\log p_1,\log a_1,\log p_2,\log a_2,\ldots\}\]
is linearly independent over $\mathbb Q$; the prime logarithms themselves have this independence by unique factorization.

For each fixed $\sigma>1$, the Euler product \eqref{eq:Q} converges absolutely and uniformly in $t=\Im s$. Kronecker's simultaneous approximation theorem allows the phases of any finite collection of $p_n^{-it}$ to approach $-1$, while those of $a_n^{-it}$ approach $+1$. Uniform convergence of the tail then gives the exact formula
\begin{equation}\label{eq:sup}
 \sup_{t\in\mathbb R}|Q_A(\sigma+it)|
 =\prod_{n\ge1}\frac{1+p_n^{-\sigma}}{1-a_n^{-\sigma}}.
\end{equation}
Applying the same argument with phases reversed gives
\begin{equation}\label{eq:invsup}
 \sup_{t\in\mathbb R}|Q_A(\sigma+it)^{-1}|
 =\prod_{n\ge1}\frac{1+a_n^{-\sigma}}{1-p_n^{-\sigma}}.
\end{equation}
Both products tend to infinity as $\sigma\downarrow1$: their logarithms dominate a positive constant times $\sum_p p^{-\sigma}$, which diverges as $\sigma\downarrow1$. Thus neither $Q_A$ nor its reciprocal is bounded on $\Re s>1$; a fortiori neither is bounded on $\Re s>\delta$ when $\delta\le1/2$.
\end{proof}

\begin{remark}[What the counterexample does and does not say]
For any \emph{fixed} $\sigma>1$, the two suprema in \eqref{eq:sup}--\eqref{eq:invsup} are finite. Unboundedness on the full half-plane arises from vertical frequencies tuned at values $\sigma\downarrow1$, not from a divergence at a fixed $s$. The perturbations can be extremely small as real numbers: the construction works inside any prescribed intervals $(p_n,p_n+\varepsilon_n)$ with $\varepsilon_n>0$. By the mean-value bound \eqref{eq:each}, $|\log Q_A(s)|\le C_K\sum_n(a_n-p_n)p_n^{-1-\sigma_K}$ on a compact set $K\subset\{\Re s>0\}$. Hence, as $\sup_n 2^n\varepsilon_n\to0$, these quotients converge to $1$ uniformly on compact sets, although each one fails to be bounded on $\Re s>1$.
\end{remark}

\begin{proof}[Proof of Theorem~\ref{thm:compare}]
The boundedness criterion is Proposition~\ref{prop:Hinfty}, the positive infinite examples are Proposition~\ref{prop:sparse}, and the negative examples are Proposition~\ref{prop:counter}.
\end{proof}

\section{Left equivalence is not similarity}
Define $C_{Q_A}$ on the logarithmic powers by $C_{Q_A}[y^s(\log y)^j]=\partial_s^j(Q_A(s)y^s)$. On $\mathcal D_k$ the diagonal comparison satisfies the formal identity
\[
 T_A^{(r)}=C_{Q_A}T^{(r)}=T^{(r)}C_{Q_A}.
\]
If $C_{Q_A}$ is boundedly invertible, it follows that
$\overline{T_A^{(r)}}=C_{Q_A}\overline{T^{(r)}}$; this is \emph{left equivalence}, not a similarity $T_A^{(r)}=C_{Q_A}T^{(r)}C_{Q_A}^{-1}$. Indeed the latter expression equals $T^{(r)}$, since the diagonal operators commute.

\begin{proposition}[No nontrivial dilation-compatible similarity]\label{prop:nosim}
Suppose a boundedly invertible operator $S$ on $S_k$ commutes with $\Delta=y\,d/dy$ on the span of its power-function eigenvectors, with the necessary domain compatibility. If it also intertwines $T^{(r)}$ with $T_A^{(r)}$ there, for either $r=2$ or $r=3$, then $Q_A\equiv1$ on $\Re s>\delta$.
\end{proposition}
\begin{proof}
The eigenspace of $\Delta$ with eigenvalue $s$ is one-dimensional, spanned by $y^s$. Hence $Sy^s=c(s)y^s$, $c(s)\ne0$. The identity $ST^{(r)}=T_A^{(r)}S$ applied to $y^s$ yields $c(s)m_r(s)=c(s)Q_A(s)m_r(s)$. Away from the discrete zero set of $m_r$, we get $Q_A(s)=1$; the identity theorem gives $Q_A\equiv1$.
\end{proof}


\section{Graph-theoretic interpretation}
\label{sec:prime-graphs}

The Euler-factor bookkeeping behind Theorem~\ref{thm:compare} has a
simple graph description, developed in full in the companion paper
\cite[Section~6]{PP26}; we record only what is needed here.

Let $G_P$ be the directed ray with vertices $p_1,p_2,\ldots$ and
edges $p_j\to p_{j+1}$.  An endpoint perturbation $a_j=p_{j+e_j}$,
$e_j\in\{0,1\}$, is a binary flow on the edges, and its signed vertex
boundary is $(\partial e)_j=e_j-e_{j-1}$ with $e_0=0$.  With
$g_s(x)=\log(1-x^{-s})$, summation by parts gives, for $\Re s>0$,
\begin{equation}\label{eq:graph-product}
 \log Q_A(s)=\sum_{n\ge1}e_n\bigl[g_s(p_n)-g_s(p_{n+1})\bigr]
 =\sum_{j\ge1}(\partial e)_j\,g_s(p_j),
\end{equation}
where the second series is summed in order.  The convergence of the
first series is the estimate used in the proof of Lemma~\ref{lem:growth}.

\begin{example}[Telescoping along the ray]\label{ex:ray}
If every edge is occupied, $\partial e=(1,0,0,\ldots)$ and
$Q_A(s)=1-2^{-s}$.  If exactly the edges $u,\ldots,v$ are occupied,
$\partial e=\mathbf 1_{\{u\}}-\mathbf 1_{\{v+1\}}$ and
$Q_A(s)=(1-p_u^{-s})/(1-p_{v+1}^{-s})$.  In both cases $Q_A$ and $1/Q_A$
are bounded on every half-plane $\Re s\ge\delta>0$, although infinitely
many, or arbitrarily many, primes move.  If instead the occupied edges
alternate, the boundary is supported at every vertex.
\end{example}

The companion paper shows that this boundary is exactly what governs
bounded comparison.  For endpoint perturbations $e,f$, the diagonal
comparison $Q_e/Q_f$ and its inverse are bounded on $\Re s>\delta$ if
and only if
\[
 \sum_{j\ge1}p_j^{-\delta}\,\bigl|(\partial e)_j-(\partial f)_j\bigr|<\infty
\]
\cite[Theorem~5.1]{PP26}.  In particular, the alternating perturbation
is not boundedly comparable with the primes for any $0<\delta\le1$.
Proposition~\ref{prop:sparse} is the special case in which the
perturbation is compared with the unperturbed primes and the changed
indices are sparse.  The counterexample of Proposition~\ref{prop:counter}
is the opposite extreme, in which no Euler factors cancel.

\section{Interpretation and next questions}
The principal distinction is between the \emph{full operator pair} and its \emph{closed joint zero-eigenspace}. Full operators need not be similar, and even the natural diagonal comparison may be unbounded. Nevertheless, the joint zero-eigenspace is invariant under all perturbations \eqref{eq:GS}; the same ambient $S_k$ norm remains available. This invariance follows from the absence of a new inner factor in $Q_A$, rather than boundedness of $Q_A$.

Research questions that remain genuinely open in this framework include:
\begin{enumerate}[leftmargin=2em]
\item Characterize arithmetic conditions on the moved primes that are both necessary and sufficient for $Q_A,1/Q_A\in H^\infty(\Re s>\delta)$ in general. The companion paper \cite{PP26} answers this for endpoint perturbations (Theorem~5.1 there) and for perturbations whose surviving Euler-factor frequencies are rationally independent (Corollary~8.1 there); the open case is that of rationally dependent interior frequencies.
\item Determine whether analytically defined $T_A^{(r)}$ have useful \emph{natural} generalized-prime dilation-series realizations, and whether those initial domains are graph cores.
\item Study perturbation behavior of resolvents, pseudospectra and operator domains \emph{outside} the common joint kernel. Invariance of the kernel by itself gives no such stability.
\item Establish how much of the theorem persists for more general Beurling systems whose Euler-product quotient has nontrivial singular inner factors or fails to lie in $N^+$.
\end{enumerate}

\noindent\textbf{Research-status qualification.} The Hardy-space and Smirnov-class mechanisms are classical. The arguments presented here are a self-contained mathematical investigation, not an assertion of publication novelty or of resolution of the Riemann hypothesis. The analytic-symbol model should not be conflated with an unconditionally convergent prime-generated dilation series.


\section{Consequences of the November 2026 RH formalization}
\label{sec:rh-2026}
On November 20, 2026 Levent Alp\"oge released a formal proof of the
Riemann hypothesis (RH), found with Anthropic's Claude \cite{Ru26}.
The Lean FRO and Mathlib Initiative certified that the Lean statement
is Mathlib's vetted statement of RH and that the machine-checked proof
is correct \cite{LF26}.  Rubinstein independently reports checks of the
original and shortened formalizations \cite{Ru26}.  Human-readable
proofs are still being developed.  The present paper does not
independently check the Lean development.\footnote{The November 2026
certification concerns RH, not the operator-theoretic assertions
studied here.  Those rest on the proofs in this paper.}

The first part of the following corollary is unconditional.  The second
part takes the certified RH theorem \cite{LF26,Ru26} as an external input.

\begin{corollary}[The joint kernel and RH]\label{cor:rh-kernel}
Let $A$ be any perturbation \eqref{eq:GS} and write
$\mathcal Z_k(A)=\Ker\overline{T_A^{(2)}}\cap\Ker\overline{T_A^{(3)}}$
in $S_k$.
\begin{enumerate}[label=(\roman*)]
\item For each $k\ge1$, $\mathcal Z_k(A)=\{0\}$ if and only if $\zeta$ has
no zeros with $\Re s>1/(2k)$.  In particular,
$\mathcal Z_1(A)=\{0\}$ if and only if RH holds.  For $k\ge2$,
$\mathcal Z_k(A)\ne\{0\}$, since $\zeta$ has zeros on the critical line.
\item Assume RH.  Then $\mathcal Z_1(A)=\{0\}$.  For $k\ge2$,
$\mathcal Z_k(A)$ is the closed $S_k$-span of
\[
 y^{\rho}(\log y)^j,\qquad \rho=\tfrac12+i\gamma,\quad\zeta(\rho)=0,
 \quad 0\le j<\operatorname{ord}_\rho\zeta,
\]
that is, of the power functions attached to \emph{all} nontrivial zeros,
and it is the same subspace for every $A$.
\end{enumerate}
\end{corollary}

\begin{proof}
By Theorem~\ref{thm:main}, $\mathcal Z_k(A)$ is the closed span of the
$y^\rho(\log y)^j$ with $\zeta(\rho)=0$ and $\Re\rho>1/(2k)$.  These
functions are nonzero, so $\mathcal Z_k(A)=\{0\}$ exactly when there is
no such $\rho$.  The nontrivial zeros lie in $0<\Re s<1$, and the
functional equation and $\zeta(\bar s)=\overline{\zeta(s)}$ make the
set of zeros symmetric under $\rho\mapsto1-\bar\rho$.  A zero with
$\Re\rho\ne\frac12$ therefore produces one with real part greater than
$\frac12$.  Hence there is no zero with $\Re s>\frac12$ if and only if
RH holds, which is the case $k=1$.  For $k\ge2$, $1/(2k)<\frac12$ and the
zeros on the critical line, for example $\rho_1$ of
Example~\ref{ex:first-zero}, give nonzero elements.  Part (ii) follows
from part (i) and Theorem~\ref{thm:main}, because under RH every
nontrivial zero has real part $\frac12$.
\end{proof}

\begin{remark}[What the corollary does and does not say]
For $k=1$ the critical line is exactly the boundary of the admissible
half-plane $\Re s>\frac12$, so no zero power lies in $S_1$.  The
statement $\mathcal Z_1(A)=\{0\}$ is the closed-operator form of the
algebraic equivalence between RH and the vanishing of the algebraic
joint kernel on $\mathcal D_1$ proved in the corrected 1981 edition.
For $k\ge2$, the dilation generator $\Delta$ has eigenvalue $\rho$ on
$y^\rho$, so the shifted generator $-i(\Delta-\frac12)$ has the real
eigenvalue $\gamma$.  It need not be symmetric or self-adjoint, and RH
does not imply that the zeros are simple.  Multiple zeros, if any,
would contribute Jordan chains $y^\rho(\log y)^j$.
\end{remark}

Grosswald--Schnitzer zero preservation places the nontrivial
zeros of each $\zeta_A$ in $\Re s>0$ on the same critical line,
with the same multiplicities and ordinates as $\zeta$.
RH neither makes the Euler quotient $Q_A$ or its reciprocal
bounded nor changes the graph-boundary summability test of
Section~\ref{sec:prime-graphs}.
It does not establish boundedly invertible similarity of the
full summation operators, equality of their nonzero spectra,
equivalence of graph domains, or pseudospectral stability.
The classical consequences of RH for the ordinary primes
(von Koch),
$\psi(x)=x+O(x^{1/2}\log^2x)$ and
$\pi(x)=\operatorname{Li}(x)+O(x^{1/2}\log x)$,
do not transfer to arbitrary Beurling or random prime systems
without additional comparison estimates.
The announced Dirichlet-$L$ GRH \cite{Ru26} is a separate result and
should not be identified with a claim for all generalized zeta functions.

\section{Glossary of technical terms}
The following definitions supplement the main text for readers
whose background is in analysis rather than generalized-prime systems.
\begin{description}[style=nextline,leftmargin=1.5em]
\item[Analytic continuation] Extension of a function beyond its original convergence domain. The continued function need not be represented by the original infinite series.
\item[Blaschke product] An inner function built from a zero sequence satisfying the Blaschke condition. Its model space records the zero set with multiplicity.
\item[Bounded diagonal comparison] A comparison quotient and its reciprocal both belong to $H^\infty$ on the chosen half-plane. It gives invertible left equivalence, not generally operator similarity.
\item[Critical line] The line $\Re s=\tfrac{1}{2}$. RH puts the nontrivial ordinary zeta zeros there.
\item[Dilation generator] The operator $\Delta=y\,d/dy$. It acts on $y^s$ as multiplication by $s$ and commutes with every dilation $f(y)\mapsto f(y/c)$.
\item[Discrete divergence] The signed incidence boundary of an edge flow. The source and target contributions telescope along prime-path blocks.
\item[Euler quotient] The function $Q_A=\zeta_A/\zeta$, an infinite product of ratios of Euler factors. For prime-gap perturbations it continues holomorphically and without zeros to $\Re s>0$.
\item[Graph core] A subspace of the domain of a closed operator that is dense for the graph norm. An identity checked on a graph core extends to the closed operator.
\item[Graph norm] The norm $\|f\|_T=(\|f\|^2+\|Tf\|^2)^{1/2}$ on the domain of an operator. It distinguishes domains even when point spectra coincide.
\item[Grosswald--Schnitzer perturbation] A sequence $a_n$ with $p_n\le a_n\le p_{n+1}$, replacing each prime by a real number in its prime gap. The zeros of the zeta function in $\Re s>0$ do not move.
\item[$H^\infty$] The algebra of bounded holomorphic functions on a half-plane, with the supremum norm. Its elements are exactly the symbols of bounded multiplication operators on $H^2$.
\item[Hardy space] A Hilbert space of analytic functions with square-integrable boundary behavior. Its multiplication operators provide the model for dilation symbols.
\item[Inner factor] A factor with boundary modulus one almost everywhere. Its Blaschke and singular components determine invariant-subspace phenomena.
\item[Joint kernel] The intersection of several operator null spaces. Pairing $r=2$ and $r=3$ removes extraneous zeros of the auxiliary factors.
\item[Kronecker approximation] A density theorem for rotations with rationally independent frequencies. It lets finitely many Euler factors approach chosen phases simultaneously.
\item[Laplace transform] The map $g\mapsto\int_0^\infty e^{-zt}g(t)\,dt$. By the Paley--Wiener theorem it identifies $L^2(0,\infty)$ with the Hardy space of the right half-plane.
\item[Operator domain] The subspace on which an unbounded operator acts. A multiplier prescription on powers is not automatically a closed operator.
\item[Outer function] A Hardy or Smirnov function determined by its boundary modulus up to phase. Outerness does not mean bounded invertibility.
\item[Pseudospectrum] The region of large resolvent norms or spectral instability under small perturbations. It is not fixed by the zeros of arithmetic symbols.
\item[Reproducing kernel] The function $K_w$ in a Hilbert space of analytic functions whose inner product with $F$ returns a multiple of $F(w)$. Power functions $y^s$ correspond to such kernels in the Hardy model.
\item[Riesz basis] The boundedly invertible image of an orthonormal basis. Having complete root vectors is weaker.
\item[Self-adjointness] Equality of an operator and its Hilbert-space adjoint, including domains. Real eigenvalues are not sufficient.
\item[Smirnov class] Ratios of bounded analytic functions with an outer denominator. This allows useful unbounded analytic quotients in the Hardy model.
\item[Sobolev threshold] The strict exponent condition $\Re s>1/(2k)$. When $k=1$ the RH critical line lies on its boundary.
\item[Spectral synthesis] Completeness of generalized eigenfunctions within the relevant invariant subspace. It requires more than matching zero sets.
\item[Weighted graph variation] A sum of weighted absolute incidence jumps. It quantifies Euler-factor cancellation for endpoint perturbations.
\item[Zero multiplicity] The vanishing order of an analytic function at a zero. Multiple zeros produce logarithmic generalized eigenvectors.
\end{description}

\section*{Acknowledgments and AI contribution}
\addcontentsline{toc}{section}{Acknowledgments and AI contribution}
The November 1, 2026 version of this paper, and the working note
\cite{JoynerNote} from which it grew, were prepared by ChatGPT (OpenAI)
for the author.  The Hardy-space model and Theorems~\ref{thm:main}
and~\ref{thm:compare}, with their proofs, originate there.  On
November 28, 2026 ChatGPT also
revised the exposition, glossary, RH status discussion, and graph
connections.  Claude (Anthropic) refereed the November 28 version and
prepared the present Draft~2.  Claude added the RH corollary
(Corollary~\ref{cor:rh-kernel}), the verification that $U_k$ is an
isomorphism, the examples, and the remark on the order of the Sobolev
norm, and replaced the duplicated graph section by a summary of the
companion paper.  Claude checked the proofs of
Lemmas~\ref{lem:growth}--\ref{lem:range}, Theorem~\ref{thm:main}, and
the propositions of Sections~\ref{sec:bounded}--\ref{sec:counter}.
AI review is not independent mathematical certification.  Details are
recorded in the separate revision log
\texttt{spectral\_invariance\_revision\_log.tex}.

\begin{thebibliography}{9}\bibitem{LF26} Lean FRO and Mathlib Initiative,
\emph{Lean verification of the Riemann Hypothesis},
joint statement, November 21, 2026.
\url{https://leanprover-community.github.io/blog/posts/riemann-hypothesis/}.
\bibitem{Ru26} M. Rubinstein,
\emph{Alp\"oge/Claude proof of the Generalized Riemann Hypothesis},
web page, November 2026, accessed November 28, 2026.
\url{https://www.math.uwaterloo.ca/~mrubinst/alpoge_claude_grh/}.

\bibitem{GS} E. Grosswald and F. J. Schnitzer, A class of modified $\zeta$ and $L$-functions, \emph{Pacific J. Math.} \textbf{74} (1978), 357--364. \href{https://msp.org/pjm/1978/74-2/pjm-v74-n2-p09-s.pdf}{journal PDF}.
\bibitem{SmirnovReference} J. S. Choa, H. O. Kim, and J. H. Shapiro, \emph{Compact composition operators on the Smirnov class}, preprint (1998), Section 2, for background on $N^+$ and inner--outer factorization. \href{https://users.math.msu.edu/users/shapiro/pubvit/Downloads/Smirnov/Smirnov.pdf}{online preprint}.
\bibitem{Fricain} E. Fricain, A. Hartmann, W. T. Ross, and D. Timotin, The Smirnov class for de Branges--Rovnyak spaces, for background on multiplier/outer-denominator methods. \href{https://pro.univ-lille.fr/fileadmin/user_upload/pages/pros/emmanuel_fricain/Smirnov-Hb-version-soumise-Result.pdf}{online preprint}.
\bibitem{PP26} D. Joyner, \emph{Operator-equivalence classes of Grosswald--Schnitzer prime perturbations}, Draft~3, November 2026, file \texttt{prime\_perturbations\_revised.tex} in this folder of \url{https://github.com/wdjoyner/alg-graph-thry}.
\bibitem{JoynerNote} ChatGPT, prepared for D. Joyner, \emph{Progress on the Spectral Questions for Zeta-Associated Dilation Operators}, November 1, 2026 (companion working note).
\bibitem{Beurling} A. Beurling, On two problems concerning linear transformations in Hilbert space, \emph{Acta Math.} \textbf{81} (1949), 239--255.
\end{thebibliography}
\end{document}